\documentclass[10pt,a4paper]{amsart}
\usepackage[margin=19mm]{geometry}
\usepackage{amsmath,amssymb,mathtools}
\usepackage[scr=rsfs,cal=euler]{mathalfa}
\usepackage{xfrac}
\usepackage[unicode,hidelinks,bookmarks=false]{hyperref}
\newcommand{\ie}{i.e.\ }
\newcommand{\cf}{cf.\ }
\newcommand{\resp}{respectively\ }
\newcommand{\Subsection}{\subsection}
\providecommand{\nobreakdash}{\nobreak\hbox{-}}
\setlength{\emergencystretch}{2em}
\multlinegap0pt
\usepackage{amssymb}
\usepackage{xcolor}
\newtheorem{theorem}{Theorem}
\def\r{\mathbb R}
\title{The Newton's problem assuming non-constant density of the fluid}
\author{Rafael L\'opez}
\date{Source: arXiv:2605.11641v1 (CC BY 4.0)}
\begin{document}
\maketitle

\noindent\emph{Source and licence.} The Newton's problem assuming non-constant density of the fluid. Version arXiv:2605.11641v1 (CC BY 4.0), distributed by arXiv under the Creative Commons Attribution 4.0 International licence, http://creativecommons.org/licenses/by/4.0/, as recorded in the arXiv licence metadata for this exact version and shown on the abstract page https://arxiv.org/abs/2605.11641v1. Full licence notice: \texttt{licenses/arxiv-2605.11641-CC-BY-4.0.txt}.\par\smallskip

\noindent\textit{Editorial scope.} Counted unit: the article's theorem t2 (source0.tex lines 202--203). The article's complete original proof of it is delivered with it (source0.tex lines 204--236, one \texttt{proof} environment of 33 lines). No mathematics is written, repaired or re-derived: the statement and the proof are the article's own lines, in the article's own notation, and no retained line is substituted or re-flowed.  Retained uncounted as the source-local context the proof needs: lines 129--131; lines 137--139; lines 192--194; lines 238--243.  One declared adaptation: exactly 1 ancillary strings inside the retained spans are omitted (image-inclusion commands and, where the source repeats a label, the repeated label definitions) and no other line differs from the source; the figure environments, with their placement, captions and labels, are kept, and the package records the omitted strings in fidelity receipts bound to the affected bodies.  No retained line contains a citation, so the wrapper carries no bibliography and no bibliography entry.  The retained lines contain the article's own diagram code, which the wrapper typesets with the standard tikz packages; no image file is delivered and the package declares no asset dependency.  Editorial formatting only, in addition: an AMS \texttt{amsart} preamble that restates the article's own declarations and macros used by the retained lines, namely the environments \texttt{theorem} on separate counters, together with the standard packages those lines need.  The retained environments print in this document as Theorem 1 in order of appearance, whereas the article prints the same items with the numbering of its own full document.  The complete native source of the article as distributed in the arXiv e-print for this exact version is bundled unchanged as \texttt{sources/200373/source0.tex}.
\begin{equation}\label{ur2}
r \left( 1 + 4 u'^2 + 3 u'^4 \right) - 2 u' \left( 1 + u'^2 \right) + 2 r u'' \left( 3 u'^2 - 1 \right) = 0.
\end{equation}

\begin{theorem}\label{t1}
The initial value problem \eqref{ur2} with initial conditions $u(0)=0$ and $u'(0)=0$ has a solution $u\in C^2([0,R])$ for some $R>0$. 
\end{theorem}

 \begin{equation}\label{ur3}
 u''=-\frac{1+4u'^2+3u'^4}{2(3u'^2-1)} +\frac{u'(1+u'^2)}{r (3u'^2-1)}.
 \end{equation}

 \begin{figure}[h]
\begin{center}
 \quad 
\end{center}
\caption{Left. The phase plane of \eqref{sys}. The orange lines are $\theta=\pm\frac{\pi}{6}$, the blue line is $\theta=\frac{\pi}{2}$. The red line is a trajectory $\alpha_\gamma$ from the point $(0,0)$ and the blue lines are different trajectories from $(1,\frac{\pi}{2}-\lambda)$ with $\lambda=0.1, 0.5$ and $0.8$.   Right. A radial solution intersecting  orthogonally the rotation axis.}\label{fig1}
\end{figure}

\begin{theorem}\label{t2} The maximal domain of a solution $u(r)$ of \eqref{ur2} with initial conditions $u(0)=u'(0)=0$ is   a bounded interval    $[0,r_M)$, where   $\lim_{r\to r_M}u'(r)=\frac{1}{\sqrt{3}}$. We also have the numerical values $r_M\approx 1.230$ and $u(r_M)\approx 0.228$.
\end{theorem}

\begin{proof}
From Thm.   \ref{t1}, we know $u''(0)=\frac14>0$, implying  that $u$ increases initially around $r=0$, and $u''$ is positive close to $0$. If $r_1>0$ were the first time that $u'$ vanished, then \eqref{ur2} would yield $u''(r_1)=\frac12$, and $r_1$ would be a minimum of $u$, which is not possible. This proves that $u'>0$ in the maximal domain of $u$. 

In order to continue  the study of the behavior of the solution $u(r)$, it is more convenient to use  parametric coordinates   $x(s)=r$ and $y(s)=u(r)$. We reparametrize   $\gamma(s)=(x(s),y(s))$ with   $x'(s)^2+y'(s)^2=1$. Then   there is and angle $\theta=\theta(s)$ such that $x'(s)=\cos\theta(s)$ and $y'(s)=\sin\theta(s)$. Equation  \eqref{ur3} is equivalent to the ODE system
\begin{equation}\label{p1}
\left\{\begin{split}
x'(s)&=\cos\theta(s),\\
y'(s)&=\sin\theta(s),\\
\theta'(s)&=\cos\theta(s)\frac{x(s)(2-\cos(2\theta(s)))-\sin(2\theta(s))}{2x(s)(2\cos(2\theta(s))-1)}.
\end{split}\right.
\end{equation}
The function $\theta'(s)$ is the curvature of the curve $\gamma(s)$.  Since  $y(s)$ does not appear in the expression for $\theta'$, and   $y(s)$ can be recovered from $\theta(s)$, we can study the behavior of the solutions $\gamma(s)$ by considering the autonomous system 
\begin{equation}\label{sys}
\left\{\begin{split}
x'&=\cos\theta,\\
\theta'&=\cos\theta\frac{x(2-\cos(2\theta))-\sin(2\theta)}{2x(2\cos(2\theta)-1)}.
\end{split}\right.
\end{equation}
 Up to $2\pi$-vertical translations in the $\theta$ variable, the phase plane of \eqref{sys} is given by 
$$\Theta=\{(x,\theta)\colon x\in\r,\theta\in(-\frac{\pi}{6},\frac{\pi}{6})\cup (\frac{\pi}{6},\pi+\frac{\pi}{6})\}.$$
 Each solution $\gamma(s)$ has associated with it unique orbit $\alpha_\gamma(s)=(x(s),\theta(s))$ of \eqref{sys}. Standard ODE theory ensures that through every point of $\Theta$ passes a unique orbit, and distinct orbits do not intersect. Let us observe that the equilibrium points are the lines $\theta= \frac{\pi}{2}+k\pi$. See Fig. \ref{fig1}. However, the flow meets these lines with a constant slope  $-\frac12$, because both equations in \eqref{sys} are multiplied by the same factor $\cos\theta$. On the other hand, and except when $x=\pm \frac{1}{\sqrt{3}}$, the flow meets  the lines $\theta=\pm\frac{\pi}{6}$ orthogonally because $\frac{d\theta}{dx}=\infty$. To analyze the behavior of the flow around the points $P_1=(\frac{1}{\sqrt{3}},\frac{\pi}{6})$ and $P_2=(-\frac{1}{\sqrt{3}},-\frac{\pi}{6})$, we   multiply both equations of \eqref{sys}  by $2x(2\cos(2\theta)-1)$, transforming them into  the autonomous system 
 
\begin{equation}\label{sys2}
\left\{
\begin{split}
 x' &= 2x(2\cos(2\theta)-1)\cos\theta, \\ 
\theta '&= \cos\theta (\sin(2\theta)+(\cos(2\theta)-2)x ).
\end{split}\right.
\end{equation}
Now, the points $P_1$ and $P_2$ are isolated equilibrium points in \eqref{sys2}. The linearization of the system at these points shows that   $P_1$ and $P_2$ are centers. Therefore, the flow around these points consists of  topological circles. As a consequence of the phase plane analysis, the flow crosses $\theta=0$ with slope $\frac12$ because $\frac{d\theta}{dx}=\frac12$ at $\theta=0$. Next, the flow moves toward the line $\theta=\frac{\pi}{6}$ being this value its limit, hence   $u'(r)\to \tan(\frac{\pi}{6})=\frac{1}{\sqrt{3}}$. Since $\theta(s)$ is increasing,  the curvature $\theta'$ of  $\gamma$ is positive in its maximal domain. This proves that the radial solution $u(r)$ is convex. 

By inspection of the phase plane, the trajectories in $\Theta$ terminate at the horizontal line $\theta=\frac{\pi}{6}$, with the value of $x(s)$ being finite in the limit. If $[0,a)$ is the maximal domain of the solutions of \eqref{p1}, then   there is $r_M<\infty$ such that $\lim_{s\to a}x(s)=r_M$.  This value $r_M$ defines the maximal domain $[0,r_M)$ of the solution $u(r)$ of \eqref{ur2} with initial conditions $u(0)=u'(0)=0$. A  graph of a solution $u(r)$ appears in   Fig. \ref{fig1}. The values of $r_M$ and $u(r_M)$ have been calculated numerically with Mathematica.
\end{proof}

\end{document}
